%
\againframe{plan}
%
\begin{frame}{Hall effect}
\begin{columns}
	\begin{column}[t]{0.8\textwidth}
		\begin{itemize}
	\item	In the presence of magnetic field acting perpendicular to the flow of electrons, the electrons experience a Lorentz force that deviates their path in the transverse direction. Accumulation of electrons along the edges leads to generation of voltage. This is called \textbf{Hall effect}.
\end{itemize}
	\end{column}
	%
	\begin{column}[t]{0.18\textwidth}
		\begin{figure}
			\includegraphics[width=\textwidth]{Edwin_Herbert_Hall__1855_1938}
			{\centering\footnotesize E. Hall }
		\end{figure}
	\end{column}
\end{columns}
%
\begin{itemize}
	\item Hall effect is prominent in two dimensional geometry and can be considered the advent of research on thin films.
	\item It was discovered before the discovery of electron!
	\item In the context of semiconductors, Hall effect has contribution from holes, in addition to electrons. It also helps determine
	\begin{itemize}
		\item the type of semiconductor (n or p),
		\item the majority carrier concentration, and
		\item the majority carrier mobility.
	\end{itemize}
\end{itemize}
\end{frame}
%
\begin{frame}{Hall effect: Lorentz force}
\begin{columns}
	\begin{column}{0.8\textwidth}
		\begin{itemize}
		\item In the presence of electric field \(\vb{E}\) and magnetic field \(\vb{B}\), the electron experiences a force given by
		\begin{align*}
			 \vb{F}_\mathrm{Lorentz} &= \vb{F}_\mathrm{elec} +  \vb{F}_\mathrm{mag} \\
			 &= -e \vb{E} -e \vb{v} \times \vb{B}
		\end{align*}
	\end{itemize}
	\end{column}
	%
	\begin{column}{0.18\textwidth}
		\begin{figure}
			\begin{overpic}[width=\textwidth]{Hendrik_Lorentz_1853_1928}
				% Overlay image on top
				\put(0,0){\includegraphics[width=0.4\textwidth]{Alfred_Nobel_1833_1896}}
			\end{overpic}
			\\
			{\centering\footnotesize H. Lorentz }
		\end{figure}
	\end{column}
\end{columns}
\end{frame}
%
\begin{frame}{Hall effect: experiment}
	\begin{figure}
		\includegraphics[width=0.32\textwidth]{hall_effect_circuit diagram_1}
		\includegraphics[width=0.32\textwidth]{hall_effect_circuit diagram_2}
		\includegraphics[width=0.32\textwidth]{hall_effect_circuit diagram_3}
	\end{figure}
	%
	\begin{itemize}
		\item The aim is to find the charge carrier type of Germanium crystal of dimensions \SI{9}{\milli\metre} x \SI{7}{\milli\metre} x \SI{0.7}{\milli\metre}.
		\item The resistivity is measured in the \textbf{four-probe geometry}. 
		\item Current is applied along \(x\) direction.
		\item Magnetic field of \SI{0.1}{\tesla} is applied along \(z\) direction. \hfill [\SI{1}{\tesla}=\SI{e4}{\gauss}]
		\item The Hall voltage is measured along \(y\) direction. 
	\end{itemize}
\end{frame}
%
\begin{frame}{Hall effect in p-type semiconductor}
\begin{columns}
	\begin{column}{0.4\textwidth}
		\begin{figure}
		\includegraphics[width=\textwidth]{hall_effect_circuit diagram_5}
	\end{figure}
	\end{column}
	%
	\begin{column}{0.58\textwidth}
			\begin{enumerate}
			\item Current along \(x\).
			\item Electron drift along \(x\).
			\item Magnetic field along \(z\).
			\item Magnetic force along \(-y\).
			\item The buildup of holes \textbf{induces} an electric field in the \(+y\) direction.
			\item The corresponding voltage across the width is positive.
		\end{enumerate}
	\end{column}
\end{columns}
%
\end{frame}
%
\begin{frame}{Hall effect in n-type semiconductor}
	\begin{columns}
		\begin{column}{0.4\textwidth}
			\begin{figure}
				\includegraphics[width=\textwidth]{hall_effect_circuit diagram_4}
			\end{figure}
		\end{column}
		%
		\begin{column}{0.58\textwidth}
			\begin{enumerate}
				\item Current along \(x\).
				\item Electron drift along \(-x\).
				\item Magnetic field along \(z\).
				\item Magnetic force along \(-y\).
				\item The buildup of electrons \textbf{induces} an electric field in the \(-y\) direction.
				\item The corresponding voltage across the width is negative.
			\end{enumerate}
		\end{column}
	\end{columns}
	%
\end{frame}
%
\begin{frame}{Hall field and Hall voltage}
	\begin{itemize}
		\item In the steady state, the magnetic field force is balanced by the induced electric field
		\[
		\vb{F} = q(\vb{E} + \vb{v} \times \vb{B}) = q(E_y - v_x B_z) = 0
		\]
		\item The induced electric field, called the \textbf{Hall field} is given by
		\begin{center}
			\framebox{$
				E_y =  v_x B_z.
			$}
		\end{center}
		\item The Hall field produces a voltage across the semiconductor, which is called the \textbf{Hall voltage} given by
		\[
		V_H = E_H w
		\]
		where \(w\) is the width of the semiconductor.
		\item The Hall voltage is positive for p-type and negative for n-type semiconductor.
	\end{itemize}
\end{frame}
%
\begin{frame}{Hall coefficient}
	\begin{definition}
		Hall coefficient \(R_H\) is defined as the Hall field per unit current density per unit magnetic field.
	\end{definition}
	%
	\begin{itemize}
		\item \(R_H\) is given by
		\[
			R_H := \frac{E_H}{j_x B_z} = \frac{v}{j_x}
		\]
		\item \(E_H\) is positive for holes and negative for electrons.
		\item However, the current density \(j\) is related to the drift velocity \(v\) by
		\[
		j = n e v \qquad [ j = p e v \quad \text{for holes}]
		\]
		\item Therefore,
		%
\begin{center}
			\framebox{$
			R_H (\mathrm{electron}) = -\frac{1}{ne} , \qquad  R_H (\mathrm{hole}) = +\frac{1}{pe}
			$}
\end{center}

	\end{itemize}
\end{frame}
%
%
\begin{frame}{Charge carrier properties}
	\begin{itemize}
		\item Semiconductor type: The sign of Hall voltage determines the type of majority carrier in a semiconductor
		\item Majority carrier density: The Hall coefficient gives the majority carrier density as
		\begin{center}
					\framebox{
				$
				n = \frac{1}{R_H e}
				$
				}
		\end{center}
		\item Majority carrier mobility: If the conductivity of the semiconductor \(\sigma\) is measured in the absence of magnetic field, then the mobility of majority carriers is given by
		\begin{center}
			\framebox{
				$			
				\mu = \frac{\sigma}{ne}
				$
				}
		\end{center}
	\end{itemize}
\end{frame}
%
\begin{frame}{Test your Understanding}
%
	\begin{problem}
		\begin{figure}
	\includegraphics[width=\textwidth]{hall_effect_test_ur_understanding}
	%\includegraphics[width=0.3\textwidth]{hall_effect_circuit diagram_5}
\end{figure}
		%
		Determine the majority carrier type, concentration and mobility. \(L = \SI{e-1}{\centi\metre}\), \(W = \SI{e-2}{\centi\metre}\),  \(d = \SI{e-3}{\centi\metre}\). \(I_x = \SI{1}{\milli\ampere}\), \(V_x = \SI{12.5}{\volt}\), \(B = \SI{500}{\gauss}\), \(V_H = \SI{-6.25}{\milli\volt}\).
	\end{problem}
\end{frame}
\begin{frame}{Applications of Hall effect}
	\begin{itemize}
		\item Magnetic sensor: Hall voltage is directly proportional to Magnetic field 
		\begin{itemize}
			\item Gaussmeter
			\item Compass sensor in smartphones
			\item Current measurement without direct contact
		\end{itemize} 
		%
		\begin{figure}
			\includegraphics[width=0.45\textwidth]{hall_effect_sensor_current}
			\includegraphics[width=0.45\textwidth]{hall_effect_sensor_rotation}
			\caption{(a) Non contact Ammeter, (b) Rotation sensor}
		\end{figure}
	\end{itemize}
\end{frame}
%
\begin{frame}{Applications of Hall effect}
	\begin{itemize}
		\item Magnetic latch (
\includegraphics[height=2ex]{kiss})	\hfill Proximity sensor
		\item Magnetic switch (
\includegraphics[height=2ex]{slap})	\hfill Rotation sensor
	\end{itemize}
	%
	\begin{figure}
		\includegraphics[width=0.9\textwidth]{hall_effect_sensor_car}
		\caption{Reference: \href{https://www.youtube.com/watch?v=tbkrJmo8Fqg}{Youtube link}}
	\end{figure}
\end{frame}
%
\begin{frame}{Summary of Hall effect}
	\begin{itemize}
		\item The Hall effect is a consequence of charged carrier moving in the presence of \textbf{perpendicular} electric and magnetic fields.
		\item The charged carrier is deflected, \textbf{inducing} a Hall voltage.
		\item The polarity of the Hall voltage depends on type of the semiconductor.
		\item The majority carrier concentration and mobility can be determined from the Hall voltage.
	\end{itemize}
\end{frame}